\documentclass[11pt,a4paper]{article}
\usepackage{amsmath,amssymb,amsthm}
\usepackage[margin=2.5cm]{geometry}
\usepackage{hyperref}
\title{Variable Cauchyriemann System}
\date{}
\begin{document}
\maketitle
\section{Definition}
The Variable Cauchy–Riemann System is an extension of the classical Cauchy-Riemann equations, adapted to variable elliptic structures. Formally, let $f = u + vi$ be a complex-valued function where $u, v \in C^1(\Omega, \mathbb{R})$ for an open domain $\Omega \subset \mathbb{R}^2$. The system is defined by the equation:

\begin{equation}
2 \partial_{\bar{z}} f = (u_x - \alpha v_y) + (v_x + u_y - \beta v_y)i + v (i_x + i_y i),
\end{equation}

where $\alpha$ and $\beta$ are functions that satisfy the ellipticity condition:

\begin{equation}
\Delta(x, y) = 4\alpha(x, y) - \beta(x, y)^2 > 0,
\end{equation}

for all $(x, y) \in \Omega$. Here, $\partial_{\bar{z}}$ is the conjugate derivative operator defined as:

\begin{equation}
\partial_{\bar{z}} = \frac{1}{2} \left( \partial_x + i \partial_y \right).
\end{equation}

This system captures the notion of holomorphicity in a variable setting, where the coefficients of the system can change with respect to the spatial variables.

\section{Examples}
Consider the variable elliptic structure defined by:

\[
\alpha(x, y) = x^2 + y^2, \quad \beta(x, y) = xy.
\]

The ellipticity condition becomes:

\[
\Delta(x, y) = 4(x^2 + y^2) - (xy)^2 > 0.
\]

To verify this, we can compute:

\[
\Delta(x, y) = 4x^2 + 4y^2 - x^2y^2.
\]

For example, at the point $(1, 1)$:

\[
\Delta(1, 1) = 4(1^2 + 1^2) - (1 \cdot 1)^2 = 8 - 1 = 7 > 0.
\]

Thus, the ellipticity condition holds. Now, if we take $f(x, y) = x + iy$, we can compute the derivatives:

\[
u_x = 1, \quad v_y = 1, \quad u_y = 0, \quad v_x = 0.
\]

Substituting these into the Variable Cauchy–Riemann System gives:

\[
2 \partial_{\bar{z}} f = (1 - \alpha) + (0 + 0 - \beta) i + (1)(i_x + i_y i).
\]

This illustrates how the variable coefficients affect the system.

\section{Key Results}
One of the significant findings regarding the Variable Cauchy–Riemann System is that covariantly holomorphic sections are not closed under pointwise multiplication. However, a scalar gauge transformation can restore closure via a natural weighted product. This result emphasizes the necessity of differentiating the imaginary unit itself, as highlighted by H. De Bie in 2012, which is crucial for maintaining internal consistency in the variable elliptic framework \cite{2601.19274v4}.

Another important result is the Rigidity-Flatness Theorem, which states that if the structure defined by $(\alpha, \beta)$ is both rigid and flat, then the coefficients $\alpha$ and $\beta$ must be constant across the domain \cite{2601.19274v4}.

\section{Diagrams / Visual Description}
The Variable Cauchy–Riemann System can be visualized in the context of a two-dimensional plane where the coefficients $\alpha(x, y)$ and $\beta(x, y)$ define a variable elliptic structure. The ellipticity condition can be represented as a region in the $(x, y)$ plane where the inequality $\Delta(x, y) > 0$ holds. This region can be shaded to indicate where the system is valid.

\begin{center}
\includegraphics[width=0.5\textwidth]{elliptic_region.png}
\caption{Region in the $(x, y)$ plane satisfying the ellipticity condition.}
\end{center}

\section{Connections}
The Variable Cauchy–Riemann System generalizes the classical Cauchy-Riemann equations and is closely related to generalized Cauchy–Riemann operators. It provides a framework for studying holomorphic functions in contexts where the underlying structure varies, such as in complex geometry and PDE theory. The system also connects to the study of variable elliptic structures, which are essential in understanding the behavior of solutions to elliptic partial differential equations.

\section{References}
\begin{itemize}
    \item \cite{2601.19274v4}
\end{itemize}\end{document}
